\documentclass[12pt,a4paper]{report}

\usepackage[a4paper,margin=1in]{geometry}
\usepackage[T1]{fontenc}
\usepackage{lmodern}
\usepackage{microtype}
\usepackage{setspace}
\usepackage{enumitem}
\usepackage{booktabs}
\usepackage{longtable}
\usepackage{array}
\usepackage{hyperref}
\usepackage{xcolor}

\hypersetup{
    colorlinks=true,
    linkcolor=black,
    urlcolor=blue,
    citecolor=black
}

\setstretch{1.15}
\setlist[itemize]{leftmargin=1.5em}
\setlist[enumerate]{leftmargin=1.8em}

\newcommand{\projecttitle}{Inventory Management System}
\newcommand{\coursename}{Database Management Systems (DBMS)}

\begin{document}

%-------------------- TITLE PAGE --------------------
\begin{titlepage}
    \centering
    \vspace*{1.5cm}
    {\Large \textbf{\coursename}\par}
    \vspace{1.5cm}
    {\Huge \bfseries Project Report Format\par}
    \vspace{0.7cm}
    {\LARGE \projecttitle\par}
    \vspace{1.8cm}

    \begin{tabular}{>{\bfseries}p{5cm} p{7cm}}
        Student Name(s): & write here \\[0.7cm]
        Roll Number(s): & write here \\[0.7cm]
        Branch \&  Section: & write here \\[0.7cm]
        Course Instructor: & Priyankar \\[0.7cm]
        Academic Session: & 2026-2027 (Monsoon Semester)
    \end{tabular}

    \vfill
    {\large Department of Computer Science}
    \vspace{0.3cm}
    {\large Indian Institute of Information Technology Bhopal}
    \vspace{1cm}
\end{titlepage}

\tableofcontents
\clearpage

%-------------------- GENERAL INSTRUCTIONS --------------------
\chapter*{General Instructions}
\addcontentsline{toc}{chapter}{General Instructions}

\begin{enumerate}
    % \item Choose a real-world application, such as Inventory Management, Library Management, Hospital Management, Student Management, or Online Shopping.
    \item The front end may be developed using any suitable programming language or framework.
    \item \textbf{MySQL is mandatory} for database creation, storage, and query execution. A different DBMS must not be used in place of MySQL.
    \item Design the database before implementing the application.
    \item Demonstrate the DBMS concepts covered in the course through the report and implementation.
    \item Include all diagrams, functional dependencies, normalization steps, SQL statements, and execution results.
    \item Submit the report, application source code, and a complete MySQL script or database export.
    \item Every team member must understand the database design and be prepared to explain it during evaluation.
\end{enumerate}

\noindent\textbf{Important:} The project title used throughout this document is an example. Students must replace it with their own approved project topic and adapt the entities, business rules, relations, and queries accordingly.

%-------------------- CHAPTER 1 --------------------
\chapter{Introduction}

\section{Project Overview}
Introduce the project and explain its purpose and intended users.

\section{Problem Statement}
Describe the problem the project aims to solve and the limitations of the existing approach.

\section{Objectives}
List the main objectives of the project.

\section{Scope of the Project}
Explain the functionalities covered by the project and its intended users.

\section{Software and Hardware Requirements}
Mention the programming language, front-end framework, MySQL version, development tools, and minimum hardware requirements.

%-------------------- CHAPTER 2 --------------------
\chapter{System Requirements and Analysis}

\section{Functional Requirements}
List the operations the system must support. For an Inventory Management System, examples include:
\begin{itemize}
    \item Adding, updating, searching, and deleting product records.
    \item Managing suppliers and product categories.
    \item Recording purchases and sales.
    \item Updating stock quantities.
    \item Generating low-stock alerts and inventory reports.
\end{itemize}


\section{Business Rules}
State the rules governing the application. Examples include:
\begin{itemize}
    \item Each product must have a unique product ID.
    \item Each supplier must have a unique supplier ID.
    \item Each sale must contain at least one product.
    \item Stock quantity must not become negative.
    \item Each purchase or sale must be associated with a valid transaction record.
\end{itemize}

%-------------------- CHAPTER 3 --------------------
\chapter{Entity--Relationship (ER) Modeling}

\section{Identification of Entities}
Identify the entities required for the project and list their attributes.

\noindent\textit{Example entities:} Product, Supplier, Category, Purchase, Sale, and Sale\_Details.

\section{Identification of Relationships}
Identify the relationships among entities and specify their cardinalities.

\noindent\textit{Examples:} A category contains multiple products; a supplier can supply multiple products; a sale can contain multiple products.

\section{ER Diagram}
Insert the complete ER diagram. It must clearly show:
\begin{itemize}
    \item Entities and their attributes.
    \item Primary keys.
    \item Relationships and relationship attributes, where applicable.
    \item Cardinality ratios (1:1, 1:N, or M:N).
    \item Total or partial participation constraints, wherever applicable.
\end{itemize}

\begin{figure}[h!]
    \centering
    \fbox{\parbox[c][4cm][c]{0.85\textwidth}{\centering Insert ER diagrams here}}
    \caption{ER diagram(s) of the proposed system}
\end{figure}

\section{Explanation of the ER Diagram}
Explain the entities, relationships, cardinalities, and participation constraints used in the diagram.

%-------------------- CHAPTER 4 --------------------
\chapter{ER-to-Relational Mapping}

\section{Mapping Rules}
Briefly state the mapping rules used to convert the ER diagram into relational tables.

\section{Step-by-Step Conversion}
Show how each entity and relationship is converted into a relation. Include the conversion of:
\begin{itemize}
    \item Strong entities.
    \item Weak entities, if applicable.
    \item 1:1 relationships.
    \item 1:N relationships.
    \item M:N relationships.
    \item Multivalued attributes, if applicable.
    \item Relationships with their own attributes, if applicable.
\end{itemize}

\section{Final Relational Schema}
Write the complete relational schema. For example:
\begin{verbatim}
PRODUCT(Product_ID, Product_Name, Price, Quantity, Category_ID)
\end{verbatim}

Underline primary keys and clearly identify foreign keys. Also identify candidate keys and composite keys, wherever applicable.

%-------------------- CHAPTER 5 --------------------
\chapter{Functional Dependencies}

\section{Identification of Functional Dependencies}
Identify the functional dependencies for each relation based on the business rules.

\noindent\textit{Example:}
\[
\text{Product\_ID} \rightarrow
\text{Product\_Name, Price, Quantity, Category\_ID}
\]

\section{Candidate Keys and Attribute Classification}
Determine the candidate keys, prime attributes, and non-prime attributes of each relation.

\section{FD Analysis}
Explain the reasoning behind the identified functional dependencies. Each dependency must be justified by the application requirements; do not list dependencies arbitrarily.

%-------------------- CHAPTER 6 --------------------
\chapter{Database Normalization}

Normalization must be demonstrated step by step, not merely stated.

\section{Unnormalized Form (UNF)}
Present an unnormalized relation, if applicable, and identify repeating groups or other problems.

\section{First Normal Form (1NF)}
Convert the relation into 1NF by ensuring that attribute values are atomic and repeating groups are eliminated.

\section{Second Normal Form (2NF)}
Identify partial functional dependencies and remove them where applicable.

\section{Third Normal Form (3NF)}
Identify transitive dependencies and decompose relations where required.

\section{Boyce--Codd Normal Form (BCNF)}
Check whether each resulting relation satisfies BCNF. If a relation violates BCNF, demonstrate the decomposition. If no violation exists, explain why the relation satisfies BCNF.

\section{Final Normalized Schema}
Present the final set of relations after normalization. Show intermediate relations and explain each decomposition. Merely stating that the database is in 3NF or BCNF is insufficient.

%-------------------- CHAPTER 7 --------------------
\chapter{Relational Integrity and Constraints}

Explain the integrity constraints applied to the database, including:
\begin{itemize}
    \item Domain constraints.
    \item Primary key constraints.
    \item Foreign key constraints.
    \item \texttt{UNIQUE} constraints.
    \item \texttt{NOT NULL} constraints.
    \item \texttt{CHECK} constraints, wherever applicable.
    \item Referential actions such as \texttt{ON DELETE} or \texttt{ON UPDATE}, wherever appropriate.
\end{itemize}

Explain how these constraints maintain data consistency and prevent invalid data.

%-------------------- CHAPTER 8 --------------------
\chapter{MySQL Database Implementation}

\section{Database Creation}
Provide the SQL commands used to create the database.

\section{Table Creation}
Provide the complete \texttt{CREATE TABLE} statements for all relations, including primary keys, foreign keys, and other constraints.

\section{Data Insertion}
Provide sample \texttt{INSERT} statements and insert sufficient meaningful data to demonstrate the system.

\section{Data Modification and Deletion}
Demonstrate appropriate \texttt{UPDATE} and \texttt{DELETE} operations.

\section{Sample Database Contents}
Show table contents using \texttt{SELECT} queries or screenshots from MySQL.

\section{SQL Query Implementation}
Write and execute queries demonstrating the following concepts:
\begin{enumerate}
    \item Simple \texttt{SELECT} queries with conditions.
    \item Sorting using \texttt{ORDER BY}.
    \item Aggregate functions: \texttt{COUNT}, \texttt{SUM}, \texttt{AVG}, \texttt{MIN}, and \texttt{MAX}.
    \item Grouping using \texttt{GROUP BY} and filtering groups using \texttt{HAVING}.
    \item Joins, including \texttt{INNER JOIN} and \texttt{LEFT JOIN}, along with other relevant joins where appropriate.
    \item Nested queries and subqueries.
    \item Queries using \texttt{IN}, \texttt{EXISTS}, \texttt{ANY}, or \texttt{ALL}, wherever relevant.
    \item Queries involving multiple tables.
    \item Queries that answer meaningful business questions.
    \item At least one view, where appropriate.
\end{enumerate}

For every query, provide the SQL statement, its purpose, and the output or a screenshot of its execution.

\section{Advanced DBMS Features}
Demonstrate at least two relevant features from the following list, preferably those covered in the course:
\begin{itemize}
    % \item Indexes and their purpose.
    \item Views.
    \item Triggers.
    % \item Stored procedures or functions.
    % \item Transactions using \texttt{START TRANSACTION}, \texttt{COMMIT}, and \texttt{ROLLBACK}.
\end{itemize}

Explain the purpose and working of each implemented feature.

%-------------------- CHAPTER 9 --------------------
\chapter{Application Development}

\section{System Architecture}
Show the interaction among the front end, application logic, and MySQL database using a simple architecture diagram.

\section{Front-End Technologies}
Mention the programming language, libraries, and frameworks used.

\section{Application Modules}
Explain the modules implemented in the project. For an Inventory Management System, examples include:
\begin{itemize}
    \item Product Management.
    \item Supplier Management.
    \item Purchase Management.
    \item Sales Management.
    \item Stock Management.
    \item Reports and Alerts.
\end{itemize}

\section{Database Connectivity}
Explain how the application connects to MySQL and executes database operations.

\section{Application Screenshots}
Include screenshots of the major application screens and their corresponding operations.

\section{Database Operations through the Application}
Demonstrate that records can be inserted, retrieved, updated, and deleted through the application wherever appropriate. Show that changes are reflected in the MySQL database.

%-------------------- CHAPTER 10 --------------------
\chapter{Testing and Results}

\section{Test Cases}
Prepare a table containing the test case ID, operation, input, expected result, actual result, and status.

\begin{center}
\small
\begin{longtable}{|p{1.1cm}|p{2.2cm}|p{2.2cm}|p{2.5cm}|p{2.5cm}|p{1.2cm}|}
\hline
\textbf{Test ID} & \textbf{Operation} & \textbf{Input} & \textbf{Expected Result} & \textbf{Actual Result} & \textbf{Status} \\
\hline
T1 & Add product & Valid product data & Product added & Fill in result & Pass/Fail \\
\hline
T2 & Duplicate ID & Existing product ID & Rejected by constraint & Fill in result & Pass/Fail \\
\hline
T3 & Invalid foreign key & Non-existent category ID & Rejected by constraint & Fill in result & Pass/Fail \\
\hline
T4 & Negative stock & Negative quantity & Rejected, if constrained & Fill in result & Pass/Fail \\
\hline
\end{longtable}
\end{center}

\section{Query Testing}
Show that the SQL queries produce the expected results.

\section{Constraint Testing}
Demonstrate that the database rejects invalid operations wherever the relevant constraints apply.

\section{Results and Discussion}
Summarize the functionality implemented and discuss the results obtained.

%-------------------- CHAPTER 11 --------------------
\chapter{Conclusion and Future Scope}

\section{Conclusion}
Summarize the project, the database concepts applied, and the main outcomes.

\section{Limitations}
Mention the limitations of the current implementation.

\section{Future Scope}
Suggest possible improvements and additional functionalities.

%-------------------- REFERENCES --------------------
\chapter*{References}
\addcontentsline{toc}{chapter}{References}
List the books, official documentation, websites, and other resources consulted during the project. Use one consistent referencing style.

%-------------------- APPENDICES --------------------
\appendix

\chapter{Complete MySQL Script}
Include the database creation commands, table definitions, constraints, sample data, SQL queries, and any views, triggers, procedures, or transactions implemented. Alternatively, provide a clearly identified supplementary SQL file with the report.

\chapter{Source Code}
Provide the source code or a repository link for the application. Ensure that the submitted code is accessible to the evaluator.

\chapter{Additional Screenshots}
Include any additional screenshots needed to demonstrate the project.

\chapter{Team Contribution}
For group projects, clearly mention the contribution of each team member.

%-------------------- MARKING SCHEME --------------------
\chapter*{Suggested Evaluation Scheme}
\addcontentsline{toc}{chapter}{Suggested Evaluation Scheme}

\begin{center}
\begin{tabular}{p{10cm}r}
\toprule
\textbf{Component} & \textbf{Marks} \\
\midrule
Problem analysis and requirements & 5 \\
ER diagram and constraints & 10 \\
ER-to-relational mapping & 10 \\
Functional dependencies and keys & 10 \\
Normalization & 15 \\
MySQL schema, constraints, and sample data & 10 \\
SQL queries and advanced DBMS features & 15 \\
Front-end integration and functionality & 10 \\
Testing and report quality & 5 \\
Viva and conceptual understanding & 10 \\
\midrule
\textbf{Total} & \textbf{100} \\
\bottomrule
\end{tabular}
\end{center}

\noindent\textbf{Evaluation principle:} Marks should primarily reflect the correctness of database design, ER-to-relational mapping, functional dependencies, normalization, SQL implementation, and understanding of DBMS concepts. Front-end appearance alone must not determine the project grade.

\end{document}
